% Native semantic reconstruction of source physical page 38 (printed page 37).
% Provenance: IEEE Std 1516-2025, rationale for Rule 7 conclusion and Rules 8--10.
\HLAStartSemanticPage{38}{37}
\fontsize{9.96pt}{11.566pt}\selectfont
\begingroup
\hbadness=10000\relax
\addfontfeatures{FakeStretch=1.035}
\vspace*{0.25pt}

By designing federates from the outset with the ability to present internal objects/attributes/interactions as\linebreak
public, the mechanisms for reuse of the federate will be in place from the start.\par
\vspace{23.52pt}

\HLARationaleHeading{rationale-8}{A.2.3 Rationale for Rule 8}
\vspace{12.20pt}

\HLANormativeRuleStatement{Federates shall be able to transfer and/or accept ownership of instance attributes dynamically during a\linebreak federation execution, as specified in their SOMs.}
\vspace{11.87pt}

By building in the functionality to transfer and accept ownership of instance attributes, federates designed\linebreak
in accordance with the HLA provide some of the basic structural tools to become federates in the widest\linebreak
possible range of future federations. With this capability, it is possible to allow a federate designed for one\linebreak
purpose to be coupled with one designed for another purpose to meet a new requirement.\par
\vspace{23.40pt}

\HLARationaleHeading{rationale-9}{A.2.4 Rationale for Rule 9}
\vspace{12.32pt}

\HLANormativeRuleStatement{Federates shall be able to vary the conditions (e.g., thresholds) under which they provide updates of\linebreak instance attributes, as specified in their SOMs.}
\vspace{11.87pt}

By designing and building in the capability to vary the conditions under which they can export instance\linebreak
attributes, federates designed in accordance with the HLA provide some of the basic structural tools to\linebreak
become federates in the widest possible range of future federations.\par
\vspace{23.40pt}

\HLARationaleHeading{rationale-10}{A.2.5 Rationale for Rule 10}
\vspace{12.32pt}

\HLANormativeRuleStatement{Federates shall be able to manage local time in a way that allows them to coordinate data exchange with\linebreak other members of a federation.}
\vspace{11.87pt}

The HLA time management structure is intended to support interoperability among federates that may rely\linebreak
on different internal concepts of time. For example, some federates might operate without any explicit\linebreak
notion of time or passage of time\HLASemicolon{} other federates may need to run faster than real-time\HLASemicolon{} and some federates,\linebreak
as is the case with some discrete event simulations, might rely upon a scrupulously coordinated causal\linebreak
order of events throughout the federation.\par
\vspace{11.92pt}

The HLA provides a single, unifying approach to time management, in that it offers a “toolbox” of services\linebreak
that support (not guarantee) federations in managing their event synchronization in a way that will meet\linebreak
both the needs of the federation and the range of internal time management strategies of the member\linebreak
federates.\par
\vspace{11.92pt}

Different categories of federates typically use a subset of an RTI’s services to implement their mechanisms.\linebreak
Federates do not explicitly indicate to the RTI how time flow is being managed (e.g., time stepped, event-\linebreak
driven, external time synchronization, independent advance). Federates for which passage of time is\linebreak
implicit may not need to use the services described in Clause 8 of IEEE Std 1516.1-2025. However, a\linebreak
federate ensure that its internal mechanisms permit coherent operation of the time management paradigm\linebreak
(or paradigms) used in the federation. Some of the considerations that occur when a federation combines\linebreak
multiple, differing time management paradigms are discussed in 8.1 of IEEE Std 1516.1-2025.\par
\endgroup
\clearpage
